\documentclass{article}
\title{株式投資の基礎 — 財務諸表と企業分析}
\author{TeX 図書館}

\begin{document}

\section{この教科書のために — 数字で会社を見る}

株価は毎秒動きますが、会社は四半期ごとにしか数字を出しません。短期の価格ノイズの裏で\textbf{ゆっくり変わる会社の価値}を読むのが株式投資の基礎体力です。この教科書では、財務三表の読み方、投資指標の意味、企業分析の型を学びます。

\section{財務三表 — 会社の 3 枚のカルテ}

\begin{center}
\begin{tabular}{|l|l|l|}
\hline
表 & 問いに答える & ポイント \\
\hline
貸借対照表(B/S) & どこに資産があり、誰のお金で? & ある\textbf{時点}の写真 \\
損益計算書(P/L) & 儲かったか? & ある\textbf{期間}の映像 \\
キャッシュフロー(CF) & お金は実際に動いたか? & 粉飾が反映しにくい \\
\hline
\end{tabular}
\end{center}

\textbf{B/S} の左側が資産(現金・在庫・設備)、右側が負債(借入)+純資産(株主の money)。\textbf{P/L} は売上から費用を引いて利益を出す階段(売上総利益 → 営業利益 → 当期純利益)。\textbf{CF} は営業・投資・財務の 3 区画に分かれます。

\begin{quote}
\textbf{【注意】}\ P/L の「利益」は会計上の判断(在庫評価・減損・引当金)が入るため、\textbf{営業キャッシュフローと突き合わせる}のが基本の検証です。「利益は増えているのに営業 CF が連年マイナス」は要注意サイン。
\end{quote}

\begin{quote}
\textbf{【演習】}\ 「売上は増加、営業 CF は減少、在庫は急増」——この組合せが示唆する可能性を 1 つ挙げよ。
\end{quote}

\textbf{解答の指針:}\ 売掛金・在庫の積み上がる「見せかけの成長」の可能性(商品が売れ残り、売上は計上されるが現金が入っていない)。

\section{PER — 利益の何倍で買っているか}

\textbf{株価収益率}( PER )は:
\[ \mathrm{PER} = \frac{\text{株価}}{\text{1 株あたり純利益(EPS)}} \]
「今期の利益の何倍の値段が付いているか」。利益の質と成長性が同じなら PER が低いほど割安。逆数 \(1/\mathrm{PER}\) は\textbf{利回り}(リターン率)として読めます——PER 12 倍なら利回り約 8.3\%。

\begin{quote}
\textbf{【注意】}\ PER が低い理由は 3 通りある:(1) 本当に割安、(2) 利益が一時的(今後落ちる)、(3) 市場が減益を織り込んでいる。\textbf{数字の裏の理由}を必ず確認するのが分析です。
\end{quote}

\begin{quote}
\textbf{【演習】}\ 株価 3000 円、EPS 150 円の会社の PER と利回り率(逆数)を求めよ。
\end{quote}

\textbf{解答の指針:}\ PER = 20 倍、リターン率 = 1/20 = 5\%。国債利回り(たとえば 1.5\%)と比較して上乗せの大きさを評価する。

\section{PBR と ROE — 純資産の視点}

\textbf{株価純資産倍率}( PBR ):
\[ \mathrm{PBR} = \frac{\text{株価}}{\text{1 株あたり純資産}} \]
1 を下回ると「解散価値より株価が安い」ことを意味します(ただし純資産の質——含み損・不要な土地——の確認は必須)。

\textbf{自己資本利益率}( ROE ):
\[ \mathrm{ROE} = \frac{\text{当期純利益}}{\text{自己資本}} \]
「株主のお金を何\%で増やしているか」の指標。\textbf{ドゥ・ポンタ分解}により 3 要素に分けられます:
\[ \mathrm{ROE} = \underbrace{\text{当期純利益率}}_{\text{儲けの質}} \times \underbrace{\text{総資産回転率}}_{\text{回す速さ}} \times \underbrace{\text{レバレッジ}}_{\text{借金の利用}} \]

\begin{quote}
\textbf{【例】}\ ROE 15\% の会社 A(利益率 20\%・回転 0.5 倍・レバ 1.5)と会社 B(利益率 2\%・回転 3 倍・レバ 2.5)は同じ ROE でも中身が全く違います。A は\textbf{高付加価値型}、B は\textbf{回転・借金型}。借金で上げた ROE は金利上昇に脆弱です。
\end{quote}

\begin{quote}
\textbf{【演習】}\ 営業利益率 10\%、総資産回転率 1.2、レバレッジ 2.0 の会社の ROE を概算せよ(簡略化のため営業利益率をそのまま使う)。
\end{quote}

\textbf{解答の指針:}\ \(0.10 \times 1.2 \times 2.0 = 24\%\)。

\section{キャッシュフローと配当 — 会社のお金の向かい先}

\textbf{フリーキャッシュフロー}( FCF )= 営業 CF − 投資 CF。「事業が生んだお金から、成長投資に再投資した残り」。会社が株主に配当・自社株買い・成長投資で振り向けられる自由な資金の源泉です。

\textbf{配当利回り} = 年配当 ÷ 株価。\textbf{配当性向} = 配当 ÷ 当期純利益(利益の何\%を配るか)。無理な配当性向の上昇は投資余力の低下と交換です。

\begin{quote}
\textbf{【演習】}\ 株価 2000 円・年配当 80 円・EPS 200 円の会社の、利回り・配当性向を求めよ。
\end{quote}

\textbf{解答の指針:}\ 利回り 4\%、配当性向 40\%。残り 60\% は内部留保(再投資)。

\section{企業分析の型 — 4 つの問い}

個別企業を見るときのチェックリスト:

\begin{enumerate}
\item \textbf{事業の質}: 何を売って稼いでいるか。10 年続く需要か。競合との差(参入障壁)は?
\item \textbf{数字の健全性}: 営業 CF と利益の整合、自己資本比率、有利子負債の水準
\item \textbf{成長の源泉}: 売上成長は量的(市場拡大)か、価格力(値上げできる力)か
\item \textbf{株主還元と価格}: PER・PBR・利回りが過去の自分や同業と比べてどうか
\end{enumerate}

\begin{quote}
\textbf{【演習】}\ 上の 4 問いを「有名コンビニチェーン」に当てはめて、それぞれの答えの調べ方を 1 行ずつ述べよ。
\end{quote}

\textbf{解答の指針:}\ 1. 決算説明資料の事業セグメントと市況。2. 有報の CF・B/S。3. 既存店売上と値上げ履歴。4. IR 資料の PER バンドや同業比較(たとえば証券アプリの指標画面)。

\section{ポートフォリオとの接続 — 個別株をどう持つか}

『ファイナンス基礎』の分散効果(§5)から、個別株の現実的な扱いが導けます:
\begin{itemize}
\item 個別株は\textbf{会社特有のリスク}(個別リスク)を持つ。分散効果は相関が低い資産間でしか働かないので、\textbf{同業他社を 5 社持っても分散にはならない}
\item インデックス(市場全体)への投資は、個別リスクをほぼ消した「市場リスクだけ」の保有
\item 個別株のリターン = 市場の動き(β)+ 会社固有の動き。固有の動きに勝つのはプロでも難しい
\end{itemize}

\begin{quote}
\textbf{【演習】}\ 「好きな会社 5 社に全資金を均等投資」の問題点を、分散効果の式の観点から 2 行で述べよ。
\end{quote}

\textbf{解答の指針:}\ 同じ市場・近い業種の 5 社は相関が高く、共分散項が残るためリスクはあまり減らない。市場全体との相関(β)もほぼそのまま残る。

\section{おわりに — 価格と価値の間}

この教科書の骨格: 財務三表(§2)で会社の実態を掴み、PER・PBR・ROE(§3・4)で値付けを測り、CF と配当(§5)で資金の流れを確かめ、4 つの問い(§6)で分析の型を持ち、ポートフォリオの位置づけ(§7)で自分の全体図に置きました。

株式投資の本質は「\textbf{価格は毎日動くが価値はゆっくり動く}」という乖離の捉え方にあります。次の一歩は『オプションの数学入門』(派生商品)、そして『個人財務の設計』(全体の生活設計)です。

\begin{quote}
\textbf{【総合演習】}
\begin{enumerate}
\item PER 15 倍・EPS 80 円の株価を求めよ。
\item ROE 12\%・配当性向 30\% の会社の「内部留保率」を求めよ。
\item 「営業 CF 黒字・投資 CF 赤字・財務 CF ほぼ 0」の会社の資金の使い方を 1 文で読み解け。
\end{enumerate}
\end{quote}

\textbf{解答の指針:}\ 1. 80 × 15 = 1200 円。2. 70\%(利益の 7 割を再投資)。3. 事業で稼いだ資金を設備・M\&A などの成長投資に投入している健全な成長段階の典型的な形。

\end{document}
